% -----------------------------------------------------------------------------
%  entries/q026.tex
%  Issued 22/10/2026. Second entry of the GKP / magic / hypergraph sequence;
%  grew out of Q025(d)(iii). All items solved before issue (pure algebra; the
%  transformation rule of (a) checked on the quarter turn F of Q025).
% -----------------------------------------------------------------------------
\begin{entry}{Q026}{No Gaussian gate goes beyond Clifford}

\entrymeta{
  topic   = {GKP code; Gaussian unitaries; symplectic group; logical Clifford group; universality},
  est     = {$\sim$2 hours including the write-up: (a) $\sim$45 min, (b) $\sim$25 min, (c) $\sim$50 min},
  origin  = Original,
  source  = {Original problem. Gaussian unitaries as affine symplectic maps:
    C.~Weedbrook \emph{et al.}, \emph{Gaussian quantum information},
    Rev.~Mod.~Phys.\ \textbf{84}, 621 (2012), Sec.~II.A.2},
  url     = https://arxiv.org/abs/1110.3234,
  prereq  = {\qref{Q025}, in particular \eqref{q025:shifts}, \eqref{q025:weyl} and
    \eqref{q025:code}},
  status  = drafted,
}

\begin{problem}
\emph{Context.} \qref{Q025} showed that the logical Fourier, sum and phase gates of the
GKP code are Gaussian unitaries. This entry proves the converse: a Gaussian unitary that
preserves the code can only ever act as a logical Clifford gate. The Clifford group is
finite up to phases, so Gaussian optics, however well controlled, cannot by itself compute
universally on GKP qudits. The missing non-Clifford gates must draw on a non-Gaussian
resource, which is where the hypergraph interactions of \qref{Q022} and magic states enter
in the following entries. Along the way you will find which symplectic maps of phase space
are compatible with the code lattice.

\emph{Setting.} Use the conventions of \qref{Q025}: one mode, the shifts $X(a)$ and $Z(b)$,
$\ell=\sqrt{2\pi/d}$, the stabilizers $S_q,S_p$ and the logical Paulis $\bar X,\bar Z$. For
$\xi=(\xi_q,\xi_p)\in\RR^2$ write
\begin{equation}
  D(\xi):=X(\xi_q)\,Z(\xi_p),\qquad D(\xi)^\dagger\,\hat r\,D(\xi)=\hat r+\xi,\qquad
  \hat r:=\begin{pmatrix}\hat q\\ \hat p\end{pmatrix},
  \label{q026:shift}
\end{equation}
so that $D(\xi)$ shifts phase space by $\xi$. A unitary $G$ is \emph{Gaussian} if it is
generated by Hamiltonians at most quadratic in $\hat q$ and $\hat p$. Equivalently (Weedbrook
\emph{et al.}, Sec.~II.A.2), it acts on the quadratures as an affine map,
\begin{equation}
  G^\dagger\hat r\,G=S\hat r+c,\qquad
  S\ \text{a real }2\times2\text{ matrix with }\det S=1,\quad c\in\RR^{2};
  \label{q026:gaussian}
\end{equation}
for one mode, $S$ is symplectic exactly when $\det S=1$. You may use that a unitary $U$ with
$U^\dagger\hat qU=\hat q$ and $U^\dagger\hat pU=\hat p$ is a multiple of the identity: the
canonical commutation relations are represented irreducibly on $L^2(\RR)$ (B.~C.~Hall,
\emph{Quantum Theory for Mathematicians}, Springer 2013, the chapter on the Stone--von
Neumann theorem). A unitary on the code space is a \emph{logical Clifford gate} if
conjugation by it maps every logical Pauli operator $\bar X^{m}\bar Z^{m'}$ to a logical
Pauli operator, up to a phase.

\begin{enumerate}[label=\textbf{(\alph*)},leftmargin=*]
  \item \emph{Gaussian unitaries move shifts.} Show that for every Gaussian unitary $G$
        and every $\xi\in\RR^{2}$,
        \begin{equation}
          G^\dagger D(\xi)\,G=e^{i\theta}\,D\bigl(S^{-1}\xi\bigr)
          \label{q026:conjugation}
        \end{equation}
        for some phase $\theta$. Check the rule on the quarter turn $F$ of \qref{Q025}.

  \item \emph{Which shifts preserve the code.} Show that $D(\xi)$ maps the code space into
        itself if and only if $\xi\in\ell\ZZ^{2}$, that is, if and only if $D(\xi)$ is, up to
        a phase, a logical Pauli operator $\bar X^{m}\bar Z^{m'}$.

  \item \emph{The lattice condition.} Let $G$ be a Gaussian unitary that maps the code
        space onto itself.
        \begin{enumerate}[label=(\roman*)]
          \item Show that $S$ has integer entries, that is, $S\in SL(2,\ZZ)$, and that $G$
                acts on the code as a logical Clifford gate, sending the labels $(m,m')$ of
                a logical Pauli operator to $S(m,m')$ modulo $d$.
          \item Read off the integer matrices of the quarter turn $F$ and of
                $e^{i\hat q^{2}/2}$ from \qref{Q025}, and decide whether a single-mode
                squeezer, $S=\mathrm{diag}(s,1/s)$ with $s>0$, can preserve the code. Is
                $S\in SL(2,\ZZ)$ also sufficient for $G$ to preserve the code? Recall what
                happened to $e^{i\hat q^{2}/2}$ for odd $d$ in \qref{Q025}.
          \item Conclude that the logical gates implementable by Gaussian unitaries form a
                finite set up to phases, and hence cannot be universal.
        \end{enumerate}
\end{enumerate}
\end{problem}

\begin{solution}[Part (a)]
\end{solution}

\begin{solution}[Part (b)]
\end{solution}

\begin{solution}[Part (c)]
\end{solution}

% \begin{remark}
% \end{remark}

\begin{takeaway}
\end{takeaway}

%  Everything below is optional.  An entry with no attempt block is complete.
%  Record only what would change what you do next time: typos and arithmetic
%  slips you caught yourself are artefacts of the medium, not facts about the
%  mathematics.  Add "public," to show this block in the public build.
% \attempt{
%   date       = ,               % ISO, YYYY-MM-DD -- required if the block is used
%   time       = ,
%   difficulty = ,               % 1-5, perceived
%   howitwent  = ,
%   insight    = ,
%   note       = ,
% }

\end{entry}